\newsection{6}{Linear Second Order Equations} {Reduction of Order}
\currentpdfbookmark{Section 5.6 Reduction of Order}{section:5.6}
\vskip14pt
\renewcommand{\thissection}{\sectiontitle{\,
REDUCTION OF ORDER}}
\thissection

\noindent
In this section we give a method for finding the general
solution of
\begin{equation} \label{eq:5.6.1}
P_0(x)y''+P_1(x)y'+P_2(x)y=F(x)
\end{equation}
if we know  a nontrivial solution $y_1$ of the complementary equation
\begin{equation} \label{eq:5.6.2}
P_0(x)y''+P_1(x)y'+P_2(x)y=0.
\end{equation}
The method is called {\color{blue}\it reduction of order\/} because it reduces
the task of solving \eqref{eq:5.6.1} to solving a first order equation.
Unlike the method of undetermined coefficients, it does not require
$P_0$, $P_1$, and $P_2$ to be constants, or $F$ to be of any special
form.

By now you shoudn't be surprised that we look for
solutions of \eqref{eq:5.6.1} in the form
\begin{equation} \label{eq:5.6.3}
y=uy_1
\end{equation}
where $u$ is to be determined so that $y$ satisfies \eqref{eq:5.6.1}.
Substituting \eqref{eq:5.6.3} and
$$\begin{array}{rcl}
y'&=& u'y_1+uy_1' \\
y''&=& u''y_1+2u'y_1'+uy_1''
\end{array}$$
into \eqref{eq:5.6.1} yields
$$
P_0(x)(u''y_1+2u'y_1'+uy_1'')+P_1(x)(u'y_1+uy_1')+P_2(x)uy_1=F(x).
$$
Collecting the coefficients of $u$, $u'$, and $u''$ yields
\begin{equation} \label{eq:5.6.4}
(P_0y_1)u''+(2P_0y_1'+P_1y_1)u'+(P_0y_1''+P_1y_1'+P_2y_1)
u=F.
\end{equation}
However, the coefficient of $u$ is zero, since $y_1$ satisfies
\eqref{eq:5.6.2}. Therefore \eqref{eq:5.6.4} reduces to
\begin{equation} \label{eq:5.6.5}
Q_0(x)u''+Q_1(x)u'=F,
\end{equation}
with
$$
Q_0=P_0y_1 \mbox{\quad and \quad}Q_1=2P_0y_1'+P_1y_1.
$$
(It isn't worthwhile to memorize the formulas for $Q_0$ and $Q_1$!)
Since \eqref{eq:5.6.5} is a linear first order equation in $u'$, we can
solve it for $u'$ by variation of parameters as in
Section~1.2, integrate the solution to obtain $u$, and then
obtain $y$ from \eqref{eq:5.6.3}.

\begin{example}\label{example:5.6.1} \rm  \mbox{} \newline
\begin{alist}
\item % (a)
 Find the general solution of
\begin{equation} \label{eq:5.6.6}
xy''-(2x+1)y'+(x+1)y=x^2,
\end{equation}
given that $y_1=e^x$ is a solution of the complementary equation
\begin{equation} \label{eq:5.6.7}
xy''-(2x+1)y'+(x+1)y=0.
\end{equation}
\item % (b)
As a byproduct of \part{a}, find a fundamental set of solutions of
\eqref{eq:5.6.7}.
\end{alist}
\end{example}

\solutionpart{a}
If
$y=ue^x$, then $y'=u'e^x+ue^x$ and
$y''=u''e^x+2u'e^x+ue^x$, so
\begin{eqnarray*}
xy''-(2x+1)y'+(x+1)y&=&x(u''e^x+2u'e^x+ue^x)\\
&&-(2x+1)(u'e^x+ue^x)+(x+1)ue^x\\
&=&(xu''-u')e^x.
\end{eqnarray*}
Therefore $y=ue^x$ is a solution of \eqref{eq:5.6.6} if and only if
$$
(xu''-u')e^x=x^2,
$$
which is a first order equation in $u'$. We rewrite it as
\begin{equation} \label{eq:5.6.8}
u''-{u'\over x}=xe^{-x}.
\end{equation}
To focus on how we apply variation of parameters to this equation, we
temporarily write $z=u'$, so that \eqref{eq:5.6.8} becomes
\begin{equation} \label{eq:5.6.9}
z'-{z\over x}=xe^{-x}.
\end{equation}
We leave it to you to show (by separation of variables) that $z_1=x$
is a solution of the complementary equation
$$
z'-{z\over x}=0
$$
for \eqref{eq:5.6.9}. By applying variation of parameters as in
Section~1.2, we can now see that every solution of
\eqref{eq:5.6.9} is of the form
$$
z=vx\mbox{\quad where \quad}
v'x=xe^{-x}, \mbox{\quad so \quad} v'=e^{-x} \mbox{\quad and \quad}
v=-e^{-x}+C_1.
$$
Since $u'=z=vx$,   $u$ is a solution of \eqref{eq:5.6.8} if and only if
$$
u'=vx=-xe^{-x}+C_1x.
$$
Integrating this yields
$$
u=(x+1)e^{-x}+{C_1\over2}x^2+C_2.
$$
Therefore the general solution of  \eqref{eq:5.6.6} is
\begin{equation} \label{eq:5.6.10}
y=ue^x=x+1+{C_1\over2}x^2e^x+C_2e^x.
\end{equation}

\solutionpart{b} By letting $C_1=C_2=0$ in \eqref{eq:5.6.10}, we see that
$y_{p_1}=x+1$ is a solution of \eqref{eq:5.6.6}. By letting $C_1=2$ and
$C_2=0$, we see that $y_{p_2}=x+1+x^2e^x$ is also a solution of
\eqref{eq:5.6.6}. Since the difference of two solutions of \eqref{eq:5.6.6} is
a solution of \eqref{eq:5.6.7},
$y_2=y_{p_1}-y_{p_2}=x^2e^x$ is a solution of \eqref{eq:5.6.7}. Since
$y_2/y_1$ is nonconstant and we already know  that $y_1=e^x$ is a
solution of \eqref{eq:5.6.6}, Theorem~\ref{thmtype:5.1.6} implies that
$\{e^x,x^2e^x\}$ is a fundamental set of solutions of
\eqref{eq:5.6.7}.\bbox

Although \eqref{eq:5.6.10} is a correct form for the general solution of
\eqref{eq:5.6.6}, it's silly to leave the arbitrary coefficient of
$x^2e^x$ as $C_1/2$ where $C_1$ is an arbitrary constant. Moreover, it's
sensible to make the subscripts of the coefficients of $y_1=e^x$ and
$y_2=x^2e^x$ consistent with the subscripts of the functions
themselves. Therefore we rewrite \eqref{eq:5.6.10} as $$
y=x+1+c_1e^x+c_2x^2e^x $$ by simply renaming the arbitrary constants.
We'll also do this in the next two examples, and in the answers to
the exercises.

\begin{example}\label{example:5.6.2}\rm  \mbox{} \newline
\begin{alist}
\item %(a)
Find the general solution of
$$
x^2y''+xy'-y=x^2+1,
$$
 given that $y_1=x$ is a solution of the complementary
equation
\begin{equation}  \label{eq:5.6.11}
x^2y''+xy'-y=0.
\end{equation}
  As a byproduct of this result, find a fundamental set of
solutions of \eqref{eq:5.6.11}.

\item %(b)
Solve the initial value problem
\begin{equation}  \label{eq:5.6.12}
x^2y''+xy'-y=x^2+1, \quad   y(1)=2,\;  y'(1)=-3.
\end{equation}
\end{alist}
\end{example}

\solutionpart{a}  If $y=ux$, then $y'=u'x+u$ and $y''=u''x+2u'$, so
\begin{eqnarray*}
x^2y''+xy'-y&=&x^2(u''x+2u')+x(u'x+u)-ux\\
&=&x^3u''+3x^2u'.
\end{eqnarray*}
Therefore $y=ux$ is a solution of \eqref{eq:5.6.12} if and only if
$$
x^3u''+3x^2u'=x^2+1,
$$
which is a first order equation in $u'$.
We rewrite it as
\begin{equation} \label{eq:5.6.13}
u''+{3\over x}u'={1\over x}+{1\over x^3}.
\end{equation}
To focus on how we apply variation of parameters to this equation, we
temporarily write $z=u'$, so that \eqref{eq:5.6.13} becomes
\begin{equation} \label{eq:5.6.14}
z'+{3\over x}z={1\over x}+{1\over x^3}.
\end{equation}
We leave it to you to show by separation of variables that
$z_1=1/x^3$ is a solution of the complementary equation
$$
z'+{3\over x}z=0
$$
for \eqref{eq:5.6.14}. By variation of parameters, every solution of
\eqref{eq:5.6.14} is of the form
$$
z={v\over x^3}\mbox{\quad where \quad}
{v'\over x^3}={1\over x}+{1\over x^3}, \mbox{\quad so \quad}
 v'=x^2+1 \mbox{\quad
and \quad} v={x^3\over 3}+x+C_1.
$$
Since $u'=z=v/x^3$,   $u$ is a solution of \eqref{eq:5.6.14} if and only
if
$$
u'={v\over x^3}={1\over3}+{1\over x^2}+{C_1\over x^3}.
$$
Integrating this yields
$$
u={x\over 3}-{1\over x}-{C_1\over2x^2}+C_2.
$$
Therefore the general solution of  \eqref{eq:5.6.12} is
\begin{equation} \label{eq:5.6.15}
y=ux={x^2\over 3}-1-{C_1\over2x}+C_2x.
\end{equation}
Reasoning as in  the solution of
Example~\ref{example:5.6.1}\part{a}, we  conclude that $y_1=x$ and $y_2=1/x$
form a fundamental set of solutions for \eqref{eq:5.6.11}.

As we explained above, we rename the constants in \eqref{eq:5.6.15} and
rewrite it as
\begin{equation}  \label{eq:5.6.16}
y={x^2\over3}-1+c_1x+{c_2\over x}.
\end{equation}

\solutionpart{b}   Differentiating  \eqref{eq:5.6.16} yields
\begin{equation} \label{eq:5.6.17}
y'={2x\over 3}+c_1-{c_2\over x^2}.
\end{equation}
 Setting $x=1$ in  \eqref{eq:5.6.16} and  \eqref{eq:5.6.17} and imposing the
initial conditions $y(1)=2$ and $y'(1)=-3$  yields
\begin{eqnarray*}
c_1+c_2&=& \phantom{-}{8\over 3} \\
c_1-c_2&=& -{11\over 3}.
\end{eqnarray*}
Solving these  equations yields $c_1=-1/2$, $c_2=19/6$.
Therefore the solution of \eqref{eq:5.6.12} is
$$
y={x^2\over 3}-1-{x\over 2}+{19\over 6x}.
$$

Using reduction of order to  find the general solution of
a homogeneous linear second order  equation leads to a
homogeneous linear first  order  equation in $u'$ that can be solved
by separation of variables. The next example illustrates this.

\begin{example}\label{example:5.6.3} \rm
 Find the general solution and a fundamental set  of solutions of
\begin{equation}  \label{eq:5.6.18}
x^2y''-3xy'+3y=0,
\end{equation}
 given that $y_1=x$ is a solution.
\end{example}

\solution
 If $y=ux$ then $y'=u'x+u$ and $y''=u''x+2u'$, so
\begin{eqnarray*}
x^2y''-3xy'+3y&=&x^2(u''x+2u')-3x(u'x+u)+3ux\\
&=&x^3u''-x^2u'.
\end{eqnarray*}
Therefore $y=ux$ is a solution of \eqref{eq:5.6.18} if and only if
$$
x^3u''-x^2u'=0.
$$
Separating the variables $u'$ and $x$ yields
$$
{u''\over u'}={1\over x},
$$
so
$$
\ln|u'|=\ln|x|+k,\mbox{\quad or, equivalently, \quad} u'=C_1x.
$$
Therefore
$$
u={C_1\over2}x^2+C_2,
$$
so the general solution of \eqref{eq:5.6.18} is
$$
y=ux={C_1\over2}x^3+C_2x,
$$
which we rewrite as
$$
y=c_1x+c_2x^3.
$$
Therefore $\{x,x^3\}$ is  a fundamental set of solutions of
\eqref{eq:5.6.18}.

\exercises
In Exercises \ref{exer:5.6.1}--\ref{exer:5.6.17} find the general solution,
given that $y_1$ satisfies the complementary equation. As a byproduct,
find a fundamental set of solutions of the complementary equation.

\begin{exerciselist}

\item\label{exer:5.6.1}
$(2x+1)y''-2y'-(2x+3)y=(2x+1)^2;  \quad y_1=e^{-x}$

\item\label{exer:5.6.2}
$x^2y''+xy'-y=\dst{4\over x^2};   \quad y_1=x$

\item\label{exer:5.6.3}
$x^2y''-xy'+y=x;   \quad y_1=x$

\item\label{exer:5.6.4}
$y''-3y'+2y=\dst{1\over1+e^{-x}};   \quad y_1=e^{2x}$

\item\label{exer:5.6.5}
$y''-2y'+y=7x^{3/2}e^x;   \quad y_1=e^x$

\item\label{exer:5.6.6}
$4x^2y''+(4x-8x^2)y'+(4x^2-4x-1)y=4x^{1/2}e^x(1+4x);   \quad
y_1=x^{1/2}e^x$

\item\label{exer:5.6.7}
$y''-2y'+2y=e^x\sec x;   \quad y_1=e^x\cos x$

\item\label{exer:5.6.8}
$y''+4xy'+(4x^2+2)y=8e^{-x(x+2)};   \quad y_1=e^{-x^2}$

\item\label{exer:5.6.9}
$x^2y''+xy'-4y=-6x-4;  \quad y_1=x^2$

\item\label{exer:5.6.10}
$x^2y''+2x(x-1)y'+(x^2-2x+2)y=x^3e^{2x};   \quad y_1=xe^{-x}$

\item\label{exer:5.6.11}
$x^2y''-x(2x-1)y'+(x^2-x-1)y=x^2e^x;  \quad y_1=xe^x$

\item\label{exer:5.6.12}
$(1-2x)y''+2y'+(2x-3)y=(1-4x+4x^2)e^x;  \quad y_1=e^x$

\item\label{exer:5.6.13}
$x^2y''-3xy'+4y=4x^4;  \quad y_1=x^2$

\item\label{exer:5.6.14}
$2xy''+(4x+1)y'+(2x+1)y=3x^{1/2}e^{-x};  \quad y_1=e^{-x}$

\item\label{exer:5.6.15}
$xy''-(2x+1)y'+(x+1)y=-e^x;  \quad y_1=e^x$

\item\label{exer:5.6.16}
$4x^2y''-4x(x+1)y'+(2x+3)y=4x^{5/2}e^{2x};  \quad y_1=x^{1/2}$

\item\label{exer:5.6.17}
$x^2y''-5xy'+8y=4x^2;  \quad y_1=x^2$

\exercisetext{In Exercises \ref{exer:5.6.18}--\ref{exer:5.6.30} find a
fundamental set of solutions, given that $y_1$ is a solution.}

\item\label{exer:5.6.18} $xy''+(2-2x)y'+(x-2)y=0;    \quad y_1=e^x$

\item\label{exer:5.6.19}
$x^2y''-4xy'+6y=0;   \quad y_1=x^2$

\item\label{exer:5.6.20}
$x^2(\ln |x|)^2y''-(2x \ln |x|)y'+(2+\ln |x|)y=0;   \quad y_1=\ln |x|$

\item\label{exer:5.6.21}
$4xy''+2y'+y=0;  \quad y_1=\sin \sqrt{x}$

\item\label{exer:5.6.22}
$xy''-(2x+2)y'+(x+2)y=0;   \quad y_1=e^x$

\item\label{exer:5.6.23}
$x^2y''-(2a-1)xy'+a^2y=0;   \quad y_1=x^a$

\item\label{exer:5.6.24}
$x^2y''-2xy'+(x^2+2)y=0;   \quad y_1=x \sin x$

\item\label{exer:5.6.25}
$xy''-(4x+1)y'+(4x+2)y=0;  \quad y_1=e^{2x}$

\item\label{exer:5.6.26}
$4x^2(\sin x)y''-4x(x\cos x+\sin x)y'+(2x\cos x+3\sin x)y=0;  \quad
y_1=x^{1/2}$

\item\label{exer:5.6.27}
$4x^2y''-4xy'+(3-16x^2)y=0;  \quad y_1=x^{1/2}e^{2x}$

\item\label{exer:5.6.28}
$(2x+1)xy''-2(2x^2-1)y'-4(x+1)y=0;  \quad y_1=1/x$

\item\label{exer:5.6.29}
$(x^2-2x)y''+(2-x^2)y'+(2x-2)y=0;  \quad y_1=e^x$

\item\label{exer:5.6.30}
$xy''-(4x+1)y'+(4x+2)y=0;  \quad y_1=e^{2x}$

\exercisetext{In Exercises~\ref{exer:5.6.31}--\ref{exer:5.6.33} solve the
initial value problem, given that $y_1$ satisfies the complementary
equation.}

\item\label{exer:5.6.31}
$x^2y''-3xy'+4y=4x^4,\quad y(-1)=7,\quad  y'(-1)=-8;   \quad y_1=x^2$

\item\label{exer:5.6.32}
$(3x-1)y''-(3x+2)y'-(6x-8)y=0, \quad   y(0)=2,\;  y'(0)=3;    \quad
y_1=e^{2x}$

\item\label{exer:5.6.33}
$(x+1)^2y''-2(x+1)y'-(x^2+2x-1)y=(x+1)^3e^x, \quad  y(0)=1,\quad y'(0)=~-1$;

$y_1=(x+1)e^x$

\exercisetext{In Exercises~\ref{exer:5.6.34} and \ref{exer:5.6.35} solve the
initial value problem and graph the solution, given that $y_1$
satisfies the complementary equation.}

\item\label{exer:5.6.34} \CGex
$x^2y''+2xy'-2y=x^2, \quad   y(1)=\dst{5\over4},\;
y'(1)=\dst{3\over2};   \quad y_1=x$

\item\label{exer:5.6.35} \CGex
$(x^2-4)y''+4xy'+2y=x+2, \quad  y(0)=-\dst{1\over3},\quad y'(0)=-1;
\quad y_1=\dst{1\over x-2}$

\item\label{exer:5.6.36}
Suppose $p_1$ and $p_2$ are continuous on $(a,b)$. Let $y_1$
be a solution of
$$
 y''+p_1(x)y'+p_2(x)y=0
\eqno{\rm (A)}
$$
that has no zeros on $(a,b)$, and let $x_0$ be in $(a,b)$.
Use reduction of order to show that $y_1$ and
$$
 y_2(x)=y_1(x)\int^x_{x_0}{1\over y^2_1(t)} \exp
\left(-\int^t_{x_0}p_1(s)\, ds\right)\,dt
$$
form a fundamental set of solutions of (A) on
$(a,b)$. \note{This exercise is related to Exercise~\ref{exer:5.1.9}.}

\color{black}
\item\label{exer:5.6.37}
The nonlinear first order equation
$$
 y'+y^2+p(x)y+q(x)=0
\eqno{\rm (A)}
$$
is  a
\href{http://http://www-history.mcs.st-and.ac.uk/Indexes/Riccati.html}
{\color{blue}\it Riccati equation\/}.  (See
Exercise~2.4.~\hspace*{-3pt}\ref{exer:2.4.55}.) Assume that $p$ and $q$ are continuous.
\begin{alist}
\item %(a)
Show that $y$ is a solution of (A) if and only if $y={z'/z}$, where
$$
 z''+p(x)z'+q(x)z=0.
\eqno{\rm (B)}
$$

\item %(b)
 Show that the general solution of (A) is
$$ y={c_1z'_1+c_2z'_2\over
c_1z_1+c_2z_2},
\eqno{\rm (C)}
$$
 where $\{z_1,z_2\}$ is a fundamental
set of solutions of (B) and $c_1$ and $c_2$ are
arbitrary constants.

\item %(c)
Does the formula  (C) imply that the first order equation
(A) has a two--parameter family of solutions?
Explain your answer.
 \end{alist}

\item\label{exer:5.6.38}
Use a method suggested by Exercise~\ref{exer:5.6.37}  to find all
solutions. of the  equation.

\begin{tabular}[t]{@{}p{168pt}@{}p{168pt}}
\part{a} $y'+y^2+k^2=0$
& \part{b} $y'+y^2-3y+2=0$
\end{tabular}

\begin{tabular}[t]{@{}p{168pt}@{}p{168pt}}
\part{c} $y'+y^2+5y-6=0$
& \part{d} $y'+y^2+8y+7=0$
\end{tabular}

\begin{tabular}[t]{@{}p{168pt}@{}p{168pt}}
\part{e} $y'+y^2+14y+50=0$
& \part{f} $6y'+6y^2-y-1=0$
\end{tabular}

\part{g} $36y'+36y^2-12y+1=0$

\item\label{exer:5.6.39}
Use a method suggested by Exercise~\ref{exer:5.6.37} and reduction of
order to find all solutions of the  equation, given that $y_1$ is
a solution.

\begin{alist}
\item %(a)
$x^2(y'+y^2)-x(x+2)y+x+2=0;  \quad y_1=1/x$

\item %(b)
$y'+y^2+4xy+4x^2+2=0;   \quad y_1=-2x$

\item %(c)
$(2x+1)(y'+y^2)-2y-(2x+3)=0;   \quad y_1=-1$

\item %(d)
$(3x-1)(y'+y^2)-(3x+2)y-6x+8=0;   \quad y_1=2$

\item %(e)
$\dst{x^2(y'+y^2)+xy+x^2- {1\over 4}=0;
\quad y_1=-\tan x -{1\over 2x}}$

\item %(f)
$\dst{x^2(y'+y^2)-7xy+7=0;  \quad y_1=1/x}$

\end{alist}

\item\label{exer:5.6.40}
The nonlinear first order equation
$$
 y'+r(x)y^2+p(x)y+q(x)=0
\eqno{\rm (A)}
$$
is  the
\href{http://http://www-history.mcs.st-and.ac.uk/Indexes/Riccati.html}
{\color{blue}\it generalized Riccati equation\/}.
(See Exercise~2.4.~\hspace*{-3pt}\ref{exer:2.4.55}.)
Assume that $p$ and $q$ are continuous and $r$ is differentiable.

\begin{alist}
\item %(a)
Show that $y$ is a solution of  (A) if and only if
 $y={z'/rz}$,  where
$$
 z''+\left[p(x)-{r'(x)\over r(x)}\right]
z'+r(x)q(x)z=0.
\eqno{\rm (B)}
$$

\item %(b)
 Show that the general solution of (A) is
$$
 y={c_1z'_1+c_2z'_2\over r(c_1z_1+c_2z_2)},
$$
where $\{z_1,z_2\}$ is a fundamental set of solutions of
(B) and $c_1$ and $c_2$ are arbitrary constants.
\end{alist}

\end{exerciselist}
